\section{Framework Description} \label{sec:framework}

This section describes the architecture and methodology used to build the
red-team vs.\ blue-team evaluation harness.

\subsection{Architecture Overview}

The harness consists of four main components connected through a shared
UDP channel and a synchronized event log:

\begin{enumerate}
    \item \textbf{Platform Simulator} (\texttt{sitl\_sim.py}): Generates
          normal MAVLink telemetry traffic (HEARTBEAT, PING,
          PARAM\_REQUEST\_LIST, STATUSTEXT) at a configurable rate,
          simulating a UAV ground-station communication link.
    \item \textbf{Attack Modules} (\texttt{attacks/}): Six standalone
          scripts, each implementing a different MAVLink injection attack.
          All share a uniform command-line interface (\texttt{--udp},
          \texttt{--rate}, \texttt{--duration}, \texttt{--out-dir}).
    \item \textbf{Detector and Defense} (\texttt{detector/}): An anomaly
          detection pipeline that trains on baseline traffic and scores
          each subsequent message.  Two defense modes are available:
          Markov-only and Isolation Forest adaptive.
    \item \textbf{Orchestrator} (\texttt{run\_demo.py}): Launches the
          simulator, starts the detector, schedules attacks (static or
          randomized), and writes synchronized logs and metrics.
\end{enumerate}

All components communicate over UDP on port 14550.  The simulator and
attack scripts send MAVLink messages to this port; the detector listens
on the same port.  Attack traffic is mixed with normal simulator traffic
at the network level, and the detector has no out-of-band knowledge of
which messages originate from attacks.

\subsection{Abstract Interfaces}

The framework defines four abstract base classes in
\texttt{framework/base.py}:

\begin{itemize}
    \item \texttt{BaseAttack}: Requires \texttt{execute(target, duration,
          rate, out\_dir)} and \texttt{stop()}.
    \item \texttt{BaseDefense}: Requires \texttt{evaluate(timestamp,
          msg\_id, src, score)} returning a block decision.
    \item \texttt{BaseDetector}: Requires \texttt{train(data)} and
          \texttt{score(observation)}.
    \item \texttt{BasePlatform}: Requires \texttt{start(target, rate)}
          and \texttt{stop()}.
\end{itemize}

New attacks or defenses can be added by subclassing the appropriate base
and registering with the orchestrator, without modifying any existing code.

\subsection{Attack Catalog}

Six MAVLink injection attacks are implemented, each targeting a different
aspect of normal traffic behavior:

\begin{table}[t]
\centering
\caption{Attack catalog with injection characteristics.}
\label{tab:attacks}
\begin{tabular}{lll}
\toprule
\textbf{Attack} & \textbf{Vector} & \textbf{Detection signal} \\
\midrule
Heartbeat flood     & Volume (high rate)       & Rate spike, distribution shift \\
Ping flood          & Volume (high rate)       & Rate spike, type ratio change \\
Param request flood & Volume (high rate)       & Rate spike, type ratio change \\
MITM identity spoof & Fake source system ID    & New source in feature window \\
Replay pattern      & Fixed message loop       & Unnatural Markov transitions \\
Command injection   & Unseen message types     & New type in feature window \\
\bottomrule
\end{tabular}
\end{table}

Each attack opens a UDP connection to port 14550, sets a spoofed source
system ID, and sends well-formed MAVLink messages at the configured rate
for the configured duration.  The detector cannot distinguish attack
traffic from simulator traffic by channel; it must rely on statistical
features.

\subsection{Feature Engineering}

The feature engine (\texttt{detector/feature\_engine.py}) maintains a
2-second sliding window over the message stream and extracts an
11-dimensional feature vector per message:

\begin{table}[t]
\centering
\caption{Feature vector components (11 dimensions).}
\label{tab:features}
\begin{tabular}{rl}
\toprule
\# & \textbf{Feature} \\
\midrule
1  & Message rate (msgs/s in window) \\
2  & Message-type entropy \\
3  & Unique source system count \\
4  & Per-source message rate \\
5  & Markov transition probability \\
6  & Inter-arrival time delta \\
7--11 & Per-type frequency ratios (HEARTBEAT, PING, \\
   & PARAM\_REQUEST, COMMAND\_LONG, STATUSTEXT) \\
\bottomrule
\end{tabular}
\end{table}

\subsection{Defense Strategies}

Two defense strategies are implemented, both operating on the same
message stream:

\subsubsection{Markov Transition Defense}
During baseline, a first-order Markov model learns transition counts
between consecutive message types.  During monitoring, the transition
probability from the previous message type to the current one is
computed.  If this probability falls below a configurable threshold
(default 0.05), the message is blocked.  This is a pattern-based defense
that relies on the assumption that attack traffic disrupts normal
message ordering.

\subsubsection{Isolation Forest Adaptive Defense}
During baseline, the Isolation Forest is trained on the 11-dimensional
feature vectors extracted from normal traffic.  During monitoring, each
message receives an anomaly score between 0 (highly anomalous) and 1
(normal).  If the score falls below a configurable threshold (default
0.3), the adaptive defense blocks the message.  This defense operates on
the full feature space and can detect deviations in rate, source
identity, message type distribution, and transition patterns
simultaneously.

\subsection{Randomized Attack Scheduling}

The orchestrator supports a \texttt{--randomize-attacks} mode that
introduces non-stationary attack traffic:
\begin{itemize}
    \item Attack order is shuffled randomly.
    \item Each attack episode is split into 5--13 segments with
          independently sampled rates and durations.
    \item Random idle gaps (benign-only traffic) are inserted between
          episodes.
    \item Short pauses are inserted between segments within an episode.
\end{itemize}
This mode tests whether defenses generalize beyond fixed attack
signatures to detect statistical deviation under varying conditions.

\subsection{Metrics}

The \texttt{ScenarioResult} class (\texttt{framework/metrics.py})
computes standardized metrics from the event log: detection rate (true
positive rate), false-positive rate, detection latency (time from attack
onset to first alert), block rate (fraction of attack traffic stopped),
accuracy recovery (post-defense accuracy minus attack-period minimum),
and mission impact (fraction of time the system is in a degraded state).
